\documentclass[../../main.tex]{subfiles}
\graphicspath{{../../image/}{../../figure/}}
\begin{document}

\section{注记}\label{section:ch08-注记}

激波的研究起源于气体动力学。Courant 与 Friedrichs 的著作 [1]（已在第 III 章引过），对过去 150 年间积累的物理与工程文献材料给出了连贯的数学阐述，为 Lax [2] 建立一般理论铺平了道路。

只要气体动力学仍是原型范式，研究的焦点就集中在严格双曲、真非线性系统上。非严格双曲系统中激波图样的复杂程度直到近年才被认识，该课题目前正处于活跃发展之中。

本章所涵盖的许多专题，其论述也可见于 Smoller [3] 与 Serre [11] 的著作。

气体动力学中 Hugoniot 轨迹的概念可追溯到 Riemann [1] 的工作；但一般设定下激波曲线的定义归功于 Lax [2]，他首先建立了\cref{theorem:ch08-激波曲线}、\cref{theorem:ch08-二阶切触}、\cref{theorem:ch08-riemann三阶} 中陈述的性质。\cref{theorem:ch08-激波曲线} 的优雅证明取自 Serre [11]。激波曲线与稀疏波曲线相重合的系统的重要性，首先由 Temple [3] 认识到，他对这些值得注意的性质作了透彻研究。详细讨论亦载于 Serre [11]。

对气体动力学，容许激波相对于其前状态应为亚声速、相对于其后状态应为超声速这一命题，见于 Riemann [1] 的开创性论文。这一原理由 Lax [2] 作为一般激波容许性准则（即 Lax $E$-条件）提出，他还证明了\cref{theorem:ch08-弱激波lax}。\cref{theorem:ch08-压缩激波存在性} 的证明见于 Ta-tsien Li 与 Wen-ci Yu [1]。另见 Hsiao 与 Chang [1]。Lax $E$-条件与稳定性之间的另一种联系由 Smoller、Temple 与 Xin [1] 建立。

真非线性缺失时的激波容许性，首先由 Bethe [1] 与 Weyl [1] 针对气体动力学系统讨论。Liu $E$-条件及相关\cref{theorem:ch08-liu蕴涵lax}、\cref{theorem:ch08-真非线性等价} 与\cref{theorem:ch08-liu蕴涵熵} 归功于 Tai-Ping Liu [2]。动机来自 Oleinik [4] 导出的 Oleinik $E$-条件，以及 Wendroff [1] 建立的 Wendroff $E$-条件。这一容许性准则显然在 1960 年代已由 Chang 与 Hsiao [1,2]（另见 Hsiao 与 Zhang [1]）预见到，但他们的工作直到很久以后才发表。

熵激波容许性条件作为连续介质热机械学基础理论的一部分已有一个多世纪（见历史引论）。对一般系统\eqref{align:ch08-守恒律一维}，形式\eqref{align:ch08-熵条件} 由 Lax [4] 提出，他建立了\cref{theorem:ch08-弱激波熵产生}、\cref{theorem:ch08-熵lax等价} 与\cref{theorem:ch08-黏性蕴涵熵}。基于方程\eqref{align:ch08-熵产生导简} 的\cref{theorem:ch08-弱激波熵产生}、\cref{theorem:ch08-接触熵零}、\cref{theorem:ch08-熵lax等价} 与\cref{theorem:ch08-liu蕴涵熵} 的证明，此处取自 Dafermos [11]。与 Liu $E$-条件等价的更严熵容许性准则版本，已由 Dafermos（见\cref{section:ch09-熵率容许性准则}）以及 Liu 与 Ruggeri [1] 提出。

条件 $2\theta-u\theta_u>0$ 与气体动力学系统\eqref{align:ch07-热弹性} 状态方程的相关性，最初由 Bethe [1] 认识到。Dafermos [30] 表明：当 Bethe 条件失效时（如发生相变的物理气体中可能出现的那样），Liu $E$-条件不再保证强激波满足热力学第二定律。

正如我们在历史引论中所见，黏性激波剖面的概念由 Rayleigh [4] 与 G.\ I.\ Taylor [1] 引入气体动力学，后由 Becker [1]、Weyl [1] 与 Gilbarg [1] 发展。系统\eqref{align:ch08-守恒律一维} 的一般形式\eqref{align:ch08-剖面ode} 首先由 Gelfand [1] 写下。

\cref{theorem:ch08-黏性等价liu} 归功于 Majda 与 Pego [1]。另见 Conlon [2]。Foy [1] 的早期论文在系统真非线性且激波为弱激波的特殊情形建立了该结果。Mock [1] 也在系统真非线性且配有凸熵的假设下证明了类似结果。刻画合适黏性矩阵 $B$ 的问题已由多位作者讨论，包括 Conley 与 Smoller [1,3]、Majda 与 Pego [1]、Serre [11]。另见 Bianchini 与 Spinolo [1]。

关于等熵（或等温）热弹性中黏性激波剖面在物理自然假设下的详细研究，见 Antman 与 Malek-Madani [1]。一般非等熵气体动力学、非凸状态方程的情形由 Pego [2] 研究，他建立了：当黏性由热导率主导时，满足 Liu $E$-条件的强激波未必允许黏性激波剖面。对未必严格双曲的一般两类守恒律系统，大振幅黏性激波剖面的存在性由 Yang、Zhang 与 Zhu [1] 讨论。系统\eqref{align:ch08-三个burgers} 及黏性矩阵\eqref{align:ch08-非恒等黏性} 由 Mailybaev 与 Marchesin [1] 详细处理。

光滑激波剖面也由其他耗散或色散机制诱导，例如：毛细作用（早期研究见 Conley 与 Smoller [2]；另见 Bertozzi 与 Shearer [1] 以及下文众多文献）；以及松弛（见 Tai-Ping Liu [21]、Yong 与 Zumbrun [1]、Dressel 与 Yong [1] 以及下文众多文献）。与 Boltzmann 方程解相关的激波剖面研究仍在进行中；见 Caflisch 与 Nicolaenko [1]，以及 Liu 与 Yu [4]。

为使问题相对简单，本章的讨论偏向严格双曲系统的中等强度激波，而对严格双曲性乃至通常双曲性失效的系统中强激波所伴随的复杂性，仅给出了提示。在很大程度上，数学物理方程中所遭遇的令人眼花缭乱的激波结构多样性，是对称性的结果——它对几何学家是喜悦之源，对分析学家是挫败之源。特别地，具有旋转不变性的系统中激波的存在性与稳定性（应用于弹性力学与磁流体力学），已由 Brio 与 Hunter [1]、Freistühler [1,2,3,4,5,6]、Freistühler 与 Liu [1]、Freistühler 与 Szmolyan [1] 讨论。关于二次或三次流函数、非严格双曲两类守恒律系统中过压缩与亚压缩激波的容许性，已有透彻研究。黏性或黏性--毛细条件，以及动力学关系，已作为容许性准则用于此目的。这一探索成败的最终检验是 Riemann 问题的适定性，将在第 IX 章讨论。从浩繁文献中取样：M.\ Shearer [2]、Schaeffer 与 Shearer [1]、Schecter 与 Shearer [1]、Schaeffer、Schecter 与 Shearer [1]、Jacobs、MacKinney 与 Shearer [1]、Schulze 与 Shearer [1]、Čanić 与 Plohr [1]、Asakura 与 Yamazaki [1]、Marchesin 与 Mailybaev [1]、Mailybaev 与 Marchesin [2,3]。其他相关文献见\cref{section:ch09-注记}。系统\cref{section:ch09-注记} 过压缩激波稳定性的讨论借自 Tai-Ping Liu [27]；另见 Tai-Ping Liu [27]。

此处未讨论的另一类跳跃间断是相界与跨声速激波，它们出现于控制相变或跨声速气流的混合型（椭圆--双曲）守恒律系统。一个原型例子是带有非单调 $\sigma(u)$ 的系统\eqref{align:ch07-p系统}，特别是经典（等熵）van der Waals 流体，其激波曲线由 LeFloch 与 Thanh [3] 描述。熵、黏性与黏性--毛细容许性准则已在该语境中试用，连同一种基于动力学关系、由微观尺度考虑驱动的新准则。关于物理背景的一部相关而全面的参考书，是 Abeyaratne 与 Knowles [3] 的有趣专著。由此类动力学关系诱导的激波通常是亚压缩的。该领域文献取样：Abeyaratne 与 Knowles [1,2]、Asakura [2]、Bedjaoui 与 LeFloch [1]、Benzoni-Gavage [4]、Fan [3,4]、Hagan 与 Slemrod [1]、Hayes 与 LeFloch [2,3]、Hayes 与 Shearer [1]、R.\ D.\ James [1]、Keyfitz [2]、Keyfitz 与 Warnecke [1]、Pego [3]、Pence [1]、Rosakis [1]、M.\ Shearer [2,5]、Slemrod [2,3,4]、Truskinovsky [1,2]。该领域由 Fan 与 Slemrod [1] 的文章以及 LeFloch [4,5] 的论文与专著充分综述。关于 Euler 方程解的容许性问题的深思熟虑的讨论，见于 Slemrod [8]。

燃烧理论中爆震波的容许性也提出有趣而困难的问题。Tai-Ping Liu 与 Tong Zhang [1] 在由\eqref{align:ch07-燃烧模型} 令 $\delta\to\infty$ 所得的高度简化的 Chapman--Jouguet 型模型语境中讨论了它。

仅在\hyperref[section:ch08-黏性激波剖面]{8.6 节}中勾勒的黏性激波剖面渐近稳定性理论，由 Serre [11] 的教科书详尽呈现。该方向的开创性论文包括：关于标量情形的 Il'in 与 Oleinik [1]；关于零过剩质量扰动系统的 Goodman [1]；引入解耦扩散波的 Tai-Ping Liu [19]；以及首次出现耦合扩散波的 Szepessy 与 Xin [1]。标量情形的权威处理见于 Serre [21] 的综述，它给出了 Freistühler 与 Serre [2,3] 的相关贡献等。关于近期用系统 Green 函数建立精确衰减率估计的综合处理，读者应参阅 Liu 与 Zeng [4] 的专著，以及同一作者的早期论文 [3]，连同 Liu 与 Yu [5]。关于接触间断、亚压缩激波、定常激波的边界效应、非严格双曲系统与各类黏性，并基于能量估计或 Green 函数的早期工作与理论推广，见 Matsumura 与 Nishihara [1,4]、Tai-Ping Liu [20,26]、Liu 与 Nishihara [1]、Liu 与 Xin [2,3]、Liu 与 Yu [1]、Tai-Ping Liu 与 Yanni Zeng [1,2]、Liu 与 Zumbrun [1]、Chern 与 Liu [1]、Lan、Liu 与 Yu [1]、Goodman、Szepessy 与 Zumbrun [1]、Kawashima 与 Matsumura [1]、Xin [3,5]、Luo 与 Xin [1]、Huang、Matsumura 与 Xin [1]、Huang Xin 与 Yang [1]、Fries [1,2]、Yanni Zeng [1,2,3]。黏性与色散在控制激波剖面渐近行为中的竞争由 Perthame 与 Ryzhik [1] 研究。配有黏性的简单模型\eqref{align:ch07-燃烧模型} 中弱爆震波的渐近稳定性由 Szepessy [3] 建立。关于阻尼效应，见 Liao、Wang 与 Yang [1]。

黏性稀疏波剖面的渐近稳定性也沿类似路线被研究；见 Liu 与 Yu [6]、Matsumura 与 Nishihara [2,3,5]、Liu、Matsumura 与 Nishihara [1]、Liu 与 Xin [1]、Xin [1,2]、Nishihara、Yang 与 Zhao [1]、Yang 与 Zhao [1]。最后，关于复合黏性波、激波和/或稀疏波的渐近稳定性，见 Huang、Li 与 Matsumura [1]，以及 Huang 与 Matsumura [1]。

关于松弛激波剖面（含接触间断）与松弛稀疏波剖面渐近稳定性的平行理论已经形成；见 Tai-Ping Liu [21]、Hsiao 与 Pan [1]、Tao Luo [1]、Luo 与 Xin [1]、Hailiang Liu [1]、Mascia 与 Natalini [1]、Mascia 与 Zumbrun [5]、Yang 与 Zhu [1]、Huijiang Zhao [2]、Ueda 与 Kawashima [1]、Feimin Huang、Ronghua Pan 与 Yi Wang [1]。类似地，Kawashima 与 Nishibata [1,2]、Kawashima 与 Tanaka [1] 建立了具有由耦合椭圆方程诱导的耗散源（辐射气体模型）的标量平衡律激波与稀疏波剖面的渐近稳定性。另见 Lattanzio、Mascia 与 Serre [1]。

激波稳定性这一在\hyperref[section:ch08-lax容许性准则]{8.3 节}中仅略微触及的课题，起源于物理文献（如 Erpenbeck [1]），但在过去三十年间已发展为一个活跃的数学研究领域。通过适当的坐标变换，激波可化为定常曲面，从而可用\cref{section:ch05-初边值问题}中概述的边值问题研究工具处理。人们可采用各种稳定性概念：谱的、线性的或非线性的，独立或互联地使用。多空间维激波的构造与稳定性理论由 Majda [2,3,4] 开创。这一开创性工作已向各方向扩展，现涵盖双曲系统的压缩、亚压缩与过压缩激波，以及改变型系统的相界。从浩繁文献中取样：Benzoni-Gavage [2,3]、Corli 与 Sablé-Tougeron [1]、Franchèteau 与 Métivier [1]、Freistühler [8]、Freistühler 与 Rohde [1]、Freistühler 与 Szmolyan [2]、Freistühler、Szmolyan 与 Wächtler [1]、Freistühler 与 Trakhinin [1]、Godin [1]、Métivier [2]、Serre [18]、Coulombel [1,2]、Zumbrun 与 Zumbrun [1]、Lewicka 与 Zumbrun [1]、Howard [1,2]、Barker、Freistühler 与 Zumbrun [1]。详尽论述见于 Benzoni-Gavage 与 Serre [2] 的专著以及 Métivier [3] 的综述文章。同一方法已用于多空间维爆震锋存在性与稳定性的研究；见 Costanzino、Jenssen、Lyng 与 Williams [1]。

与上述在精神与技术上密切相关的，是一维黏性、毛细、辐射与松弛激波剖面、相界与爆震锋的稳定性理论。这是近年来活跃的研究领域，并仍在迅速发展。Zumbrun [5,6,10] 的综述文章给出了详尽论述及 exhaustive 的参考文献列表。相关贡献取样：Zumbrun [3,4,7,8,9]、Gardner 与 Zumbrun [1]、Howard 与 Zumbrun [1,2]、Benzoni-Gavage、Serre 与 Zumbrun [1,2]、Benzoni-Gavage [8]、Guès、Métivier、Williams 与 Zumbrun [1,2,3,4,7]、Hoff 与 Zumbrun [1,2]、Lyng 与 Zumbrun [1]、Plaza 与 Zumbrun [1]、Mascia 与 Zumbrun [1,2,3,4,5,6]、Texier 与 Zumbrun [1]、Barker、Humpherys、Rudd 与 Zumbrun [1]、Howard、Raoofi 与 Zumbrun [1]、Métivier、Texier 与 Zumbrun [1]、Howard [1]、Jenssen、Lyng 与 Williams [1]、Barker、Lewicka 与 Zumbrun [1]、Nguyen 与 Zumbrun [1]、Ha 与 Yu [1]、Härterich [2]、Kreiss、Kreiss 与 Lorenz [1]、Luo、Rauch、Xie 与 Xin [1]。类似思想与技术适用于带黏性耗散的守恒律系统周期解的稳定性研究；见 Oh 与 Zumbrun [1,2]、Beck、Sandstede 与 Zumbrun [1]、Freistühler 与 Szmolyan [3]、Serre [23]。

关于由双曲守恒律离散化所得有限差分系统的激波剖面稳定性，一门平行理论正在形成；见 Jian-Guo Liu 与 Zhou Ping Xin [1,2]、Jiang 与 Yu [1]、Liu 与 Yu [2]、Benzoni-Gavage [6]。

\hyperref[section:ch08-非守恒激波]{8.7 节}中勾勒的、非守恒形式拟线性双曲系统弱解的概念，由 LeFloch [2]、Dal Maso、LeFloch 与 Murat [1] 引入。关于这些思想的发展与应用，见 Amadori、Baiti、LeFloch 与 Piccoli [1]、Hayes 与 LeFloch [2,3]、LeFloch 与 Tzavaras [1,2]，以及\cref{section:ch14-注}中的文献。关于各种应用的综述与列表，见 LeFloch [4]，以及 Berthon、Coquel 与 LeFloch [1]。关于稳定性问题，见 Xiao-Biao Lin [1]。

\end{document}
